% =====================================================================
% Sostituisce §V-A-2 (calibrazione, protocollo di adattamento).
% Numeri: benchmark_1s1t_iwcv_summary.csv, benchmark_1s1t_prior_summary.csv,
% sens_tuned_worlds_1src_{observational,concept}_shift.csv (2026-08-21).
% =====================================================================

\subsubsection{Calibration}
Calibration responds to shift in a way discrimination does not. Under
label-prior shift, prior adjustment reduces target ECE from $0.217$ to $0.088$
against an oracle floor of $0.060$, and this is the one pairing in which a
method's assumption matches the mechanism the generator applied. The direction
was stated before the runs: the same method degrades calibration under every
other regime, and it does so in $66$ of $100$ configurations under covariate
shift and $70$ of $100$ under measurement shift.

Two adversarial methods also improve calibration, DANN under measurement shift
and CDAN under concept shift. Both effects hold across the configuration space
rather than at the benchmark point alone. Over $100$ sampled generator
configurations DANN lowers target ECE in $81$ of them under measurement shift
and CDAN in $86$ under concept shift. The size varies: where the interval
excludes zero the median reduction is $0.033$ and $0.037$, and where it does not
the median is $0.003$ and $0.004$. The direction is stable throughout, and the
effect falls below what ten seeds resolve in roughly half the space.

The benefit depends on selecting hyperparameters for each run. With the
hyperparameters held at a single setting, the same two methods raise target ECE
instead of lowering it, in $75$ and $93$ configurations of $100$. Both patterns
are far from an even split, so the reversal is a property of the selection and
not sampling noise. A practitioner who applies these methods at default settings
should expect the opposite of the published effect.

Neither effect is specific to the mechanism it accompanies. Under no shift at
all, DANN lowers target ECE by $0.010$ and CDAN by $0.013$, which is most of
what they achieve under the regimes where they help. Adversarial training
produces better calibrated probabilities whether or not there is a shift to
correct, so an observed calibration gain carries little information about which
mechanism is present. Prior adjustment is the exception: its gain of $0.136$
over the no-shift level appears under label-prior shift and nowhere else.

Platt scaling improves calibration more broadly and assumes nothing about the
mechanism, reducing ECE in $66$ of $100$ configurations under measurement shift
and $73$ of $100$ under concept shift. No single strategy improves
discrimination and calibration together.
